\chapter{CertOS Artifact Audit}
\label{app:certos-artifact-audit}

This appendix records the bounded runtime evidence promoted into
Chapter~\ref{ch:certos-runtime-assurance}.  The runtime chapter itself carries the
main story.  The present appendix preserves the implementation-facing tables
that a committee or reviewer may reasonably ask for once the runtime is part
of the defense package.

\section{Artifact Bundle}

\begin{table}[ht]
\centering
\caption{CertOS artifact bundle and its role in the thesis.}
\label{tab:certos-artifact-bundle}
\begin{tabular}{p{4.8cm}p{2.6cm}p{4.0cm}}
\toprule
\textbf{Artifact} & \textbf{Type} & \textbf{Role} \\
\midrule
\path{reports/certos/certos_lifecycle.csv} & CSV & Stepwise battery runtime trace \\
\path{reports/certos/certos_summary.json} & JSON & Aggregate counts by lifecycle phase \\
\path{reports/certos/latency_report.csv} & CSV & Runtime latency and throughput \\
\path{reports/certos/audit_completeness.json} & JSON & Audit-chain completeness check \\
\path{reports/certos/recovery_report.md} & Markdown & Recovery note after fallback \\
\path{reports/certos/certos_vehicle_toy.json} & JSON & Bounded non-battery lifecycle toy \\
\path{reports/certos/invariant_tests.log} & Log & Runtime invariant test surface \\
\path{reports/certos/audit_ops.jsonl} & JSONL & Append-only lifecycle operation stream \\
\bottomrule
\end{tabular}
\end{table}

\section{Lifecycle Counts by Status}

\begin{table}[ht]
\centering
\caption{Status counts derived from the locked CertOS summary artifact.}
\label{tab:certos-status-counts}
\begin{tabular}{lr}
\toprule
\textbf{Status bucket} & \textbf{Count} \\
\midrule
valid & 47 \\
degraded & 16 \\
fallback & 33 \\
total steps & 96 \\
audit entries & 96 \\
runtime interventions & 33 \\
\bottomrule
\end{tabular}
\end{table}

The lifecycle counts matter because they make the runtime story concrete.
The battery runtime spends roughly half of the governed demo in ordinary
\texttt{valid} operation, moves through a smaller degraded regime when the
certificate horizon narrows, and explicitly falls back when the horizon is
exhausted.  That is the empirical meaning of the runtime chapter's bounded
claim.

\section{Certificate Field Schema}

The certificate engine stores only a small number of fields, but those
fields are enough to reconstruct the runtime decision trail.  The runtime
then augments the certificate with a previous-hash link and a current
certificate hash.

\begin{table}[ht]
\centering
\caption{Certificate and runtime-state fields used by CertOS.}
\label{tab:certos-certificate-schema}
\begin{tabular}{p{3.0cm}p{1.8cm}p{5.8cm}}
\toprule
\textbf{Field} & \textbf{Type} & \textbf{Meaning} \\
\midrule
\texttt{op} & string & Lifecycle operation that created the certificate \\
\texttt{proposed\_action} & mapping & Action before runtime filtering \\
\texttt{safe\_action} & mapping & Action released after runtime filtering \\
\texttt{validity\_horizon} & integer & Remaining certified horizon \\
\texttt{status} & string & \texttt{valid}, \texttt{degraded}, \texttt{expired}, or \texttt{revoked} \\
\texttt{prev\_hash} & string / null & Previous certificate hash in the ledger chain \\
\texttt{cert\_hash} & string & Current certificate hash \\
\texttt{step} & integer & Runtime step index in the governed demo \\
\texttt{fallback\_active} & boolean & Whether fallback action is currently active \\
\texttt{audit\_count} & integer & Number of audit entries emitted so far \\
\bottomrule
\end{tabular}
\end{table}

\section{Audit Ledger Record Format}

\begin{table}[ht]
\centering
\caption{Audit ledger record fields from the CertOS audit-ledger module.}
\label{tab:certos-audit-ledger-schema}
\begin{tabular}{p{2.6cm}p{1.7cm}p{6.1cm}}
\toprule
\textbf{Field} & \textbf{Type} & \textbf{Meaning} \\
\midrule
\texttt{op} & string & Lifecycle operation appended to the ledger \\
\texttt{step} & integer & Runtime step where the event occurred \\
\texttt{cert\_hash} & string / null & Hash anchor for the relevant certificate \\
\texttt{meta} & mapping & Optional metadata such as recovery details \\
\bottomrule
\end{tabular}
\end{table}

The ledger is intentionally austere.  It is not a full event-sourcing
platform.  It is an append-only trail sufficient to prove that lifecycle
events happened, in order, and with the expected linkage to certificate
records.

\section{Latency and Recovery Audit}

\begin{table}[ht]
\centering
\caption{Latency and recovery audit from the CertOS runtime artifacts.}
\label{tab:certos-latency-recovery}
\begin{tabular}{p{5.8cm}r}
\toprule
\textbf{Metric} & \textbf{Value} \\
\midrule
Runtime step latency (ms) & 0.59 \\
Runtime throughput (steps/s) & 1691.70 \\
Audit completeness (\%) & 100.0 \\
Recovery transitions noted & yes \\
\bottomrule
\end{tabular}
\end{table}

The recovery note is qualitative rather than mathematical: the runtime
artifact records that transitions from fallback back to valid execution were
observed.  This is enough for a bounded systems chapter.  It is not enough
to claim optimal recovery or universal recovery guarantees.

\section{Bounded Second-Domain Runtime Surface}

\begin{table}[ht]
\centering
\caption{Vehicle-shaped CertOS toy artifact.}
\label{tab:certos-toy-domain}
\begin{tabular}{lr}
\toprule
\textbf{Metric} & \textbf{Value} \\
\midrule
Domain & vehicle \\
Steps & 12 \\
Fallback steps & 4 \\
Audit entries & 12 \\
Interventions & 4 \\
\bottomrule
\end{tabular}
\end{table}

This toy artifact earns its place in the appendix because it helps defend a
very narrow point: the runtime lifecycle is adapter-shaped rather than
battery-field hard coded.  It does not widen the thesis proof boundary.
